% =========================================================
% Proof: GCD Remainder Lemma
% Source: volume-iii/algebra/abstract-algebra/notes/notes-integers.tex
% =========================================================

\subsection*{GCD Remainder Lemma}
\label{prf:gcd-remainder}

\begin{remark*}[Return]
\hyperref[lem:gcd-remainder]{$\leftarrow$ Back to Lemma (GCD Remainder Lemma) in Notes}
\end{remark*}

\begin{proof}
\Claim Let $a, b, q, r \in \mathbb{Z}$ with $b \neq 0$ and $a = bq + r$.
Then $\gcd(a,b) = \gcd(b,r)$.

\Let $d = \gcd(a,b)$ and $d' = \gcd(b,r)$.
\Strategy We show $d = d'$ by mutual divisibility.

\medskip
\noindent\textit{$d \mid d'$.}

\ByThm{Bézout's Identity}, $d = ma + nb$ for some $m, n \in \mathbb{Z}$.
Substituting $a = bq + r$:
\[
    d = m(bq + r) + nb = b(mq + n) + mr.
\]
So $d$ is an integer linear combination of $b$ and $r$, hence a common
divisor of $b$ and $r$.
\ByDef{Greatest Common Divisor}, since $d'$ is the greatest common divisor
of $b$ and $r$, we have $d \mid d'$.

\medskip
\noindent\textit{$d' \mid d$.}

\ByThm{Bézout's Identity}, $d' = m'b + n'r$ for some $m', n' \in \mathbb{Z}$.
Substituting $r = a - bq$:
\[
    d' = m'b + n'(a - bq) = n'a + b(m' - n'q).
\]
So $d'$ is an integer linear combination of $a$ and $b$, hence a common
divisor of $a$ and $b$.
\ByDef{Greatest Common Divisor}, since $d$ is the greatest common divisor
of $a$ and $b$, we have $d' \mid d$.

\medskip
Since $d \mid d'$ and $d' \mid d$ with $d, d' > 0$, we conclude $d = d'$.

\Hence $\gcd(a,b) = \gcd(b,r)$. \AsReq
\end{proof}

\begin{remark*}[Proof shape]
Mutual divisibility. Both directions follow the same pattern: apply
Bézout's Identity to express one gcd as a linear combination, substitute
the relation $a = bq + r$ (or $r = a - bq$) to rewrite it as a linear
combination of the other pair, then invoke the universal-divisor property
of the gcd. The substitution is the only real work — the rest is
bookkeeping.

\FlashProof{Mutual divisibility. Both directions follow the same pattern: apply Bézout's Identity to express one gcd as a linear combination, substitute the relation $a = bq + r$ (or $r = a - bq$) to rewrite it as a linear combination of the other pair, then invoke the universal-divisor property of the gcd. The substitution is the only real work — the rest is bookkeeping.}
\FlashUses{\begin{itemize} \item \hyperref[thm:bezout-identity]{Bézout's Identity} — to express each gcd as an integer linear combination of its arguments. \item \hyperref[def:gcd]{Definition (Greatest Common Divisor)} — the universal-divisor property: any common divisor divides the gcd. \end{itemize}}
\end{remark*}

\begin{remark*}[Dependencies]
\begin{itemize}
    \item \hyperref[thm:bezout-identity]{Bézout's Identity} — to express
          each gcd as an integer linear combination of its arguments.
    \item \hyperref[def:gcd]{Definition (Greatest Common Divisor)} —
          the universal-divisor property: any common divisor divides
          the gcd.
\end{itemize}
\end{remark*}